\chapter{Supervised Machine Learning}

\section{Why machine learning in an econometrics language?}

Classical econometrics excels when the goal is inference: estimating a
parameter, testing a hypothesis, or recovering a causal effect. Machine
learning is most useful when the goal is prediction: finding a flexible
function $\hat f(X)$ that approximates $Y$ well out of sample.

\hayashi{} bridges the two. Its ML commands use the same formula syntax as
OLS or logit, and they return model objects that work with \texttt{predict},
\texttt{glance}, and \texttt{tidy}.

This chapter covers tree-based methods, neural networks, and a Gaussian
process. It is not a full machine-learning course; it shows how to
estimate, interpret, and compare these models inside \hayashi{}.

\section{Random Forest}

Random Forest (Breiman, 2001) fits an ensemble of decision trees on
bootstrap samples, averaging their predictions to reduce variance.

\begin{lstlisting}[caption={Random forest regression}]
set_seed(42)
let n = 500
generate df x1 = rnormal(n)
generate df x2 = rnormal(n)
generate df y = 2 + 0.5*x1 - 0.3*x2 + 0.4*x1*x2 + rnormal(n)

let m_rf = rf(y ~ x1 + x2, df, trees=200, depth=10)
print(m_rf)

let g = glance(m_rf)
print(g)
\end{lstlisting}

Key outputs:

\begin{itemize}
\item \texttt{r2} and \texttt{oob\_r2}: in-sample and out-of-bag $R^2$
\item feature importance: reduction in impurity attributable to each regressor
\end{itemize}

Use \texttt{predict(df, yhat = m\_rf)} to obtain fitted values on a DataFrame.

\section{Gradient Boosting and XGBoost}

Gradient Boosting builds trees sequentially, each one fitting the
residuals of the previous ensemble. XGBoost (Chen and Guestrin, 2016)
extends this with second-order information and regularization.

\begin{lstlisting}[caption={XGBoost regression}]
let m_xgb = xgboost(y ~ x1 + x2, df, trees=100, depth=4, lr=0.1)
print(glance(m_xgb))
\end{lstlisting}

Boosting tends to outperform Random Forest on structured tabular data,
but it is more sensitive to hyper-parameters and can overfit.

\section{Multilayer Perceptron}

The \texttt{mlp} command fits a single-hidden-layer feed-forward network with
sigmoid activation and back-propagation with momentum.

\begin{lstlisting}[caption={Neural network}]
let m_mlp = mlp(y ~ x1 + x2, df, hidden=20, lr=0.01, epochs=500)
print(glance(m_mlp))
\end{lstlisting}

It is useful for uncovering nonlinear interactions, but it offers less
transparency than trees and often needs tuning.

\section{LSTM and Transformer for series}

For time-series, \texttt{lstm} and \texttt{transformer} use sequence windows to
predict the next value. They standardise the target internally.

\begin{lstlisting}[caption={LSTM forecast}]
let n = 200
generate df y = 0.1*(_n - n/2) + 0.5*sin(2*pi()*_n/20) + rnormal(n)
let m_lstm = lstm(df, y, seqlen=10, hidden=10, epochs=100, forecast=5)
print(m_lstm)
\end{lstlisting}

These models are deterministic once the seed is set, but their
predictive quality is convention-sensitive. They are experimental; to use
\texttt{lstm} and \texttt{transformer} build Hayashi with \texttt{--features experimental}.

\section{Gaussian Process}

A Gaussian Process defines a distribution over functions. It is
non-parametric in the sense that it does not fix a functional form, and
it reports uncertainty around each prediction.

\begin{lstlisting}[caption={Gaussian process regression}]
let m_gp = gp(y ~ x1 + x2, df)
print(m_gp)
\end{lstlisting}

GP is attractive for small datasets and when uncertainty propagation
matters. It is slower and more memory-intensive than tree methods.

\section{When to prefer which?}

\begin{itemize}
\item \textbf{Interpretability + robustness:} Random Forest or Gradient Boosting
\item \textbf{Predictive accuracy with tuning:} XGBoost
\item \textbf{Nonlinear smooth functions:} MLP or GP
\item \textbf{Sequences:} LSTM / Transformer
\item \textbf{Uncertainty:} GP or BART
\end{itemize}

All commands accept a standard Hayashi formula and return an object
compatible with \texttt{predict} and \texttt{glance}. Standard errors are generally not
reported because these methods are not designed for direct inference on
individual coefficients.
